% -----------------------------------------------------------------------------
% Tangent-centred, unprojected Hermite--ONS.
%
% This patch is written for insertion into the manuscript
% "Direct Hermite Projections for Stochastic Bandit Convex Optimization".
% It uses the manuscript's notation, filtration, extension e, Hermite
% coefficients g_t,H_t, regret-switching lemma, bracket proposition, trace
% lemma, Freedman inequality, self-normalised martingale inequality, and
% chi-square tail lemma.
%
% The theorem below is unconditional and changes no algorithmic operation
% beyond (i) deleting the projection in the mean update and (ii) replacing the
% raw Hermite estimators by the tangent-centred estimators.  The theorem uses a
% smaller value of eta.  This yields coarse, automatically proved confinement
% and regret O(d^3 M sqrt(n) ell^4).  The final subsection records the exact
% obstruction to recovering the original eta by the most direct stopped
% distance-to-K_epsilon argument.
% -----------------------------------------------------------------------------

\subsection{Tangent-centred unprojected Hermite--ONS}
\label{sec:ctr-unprojected}

Write
\begin{align}
    p(x)=\frac{\pi(x)}{1-\eps},
    \qquad
    \K_\eps=\{x\in\R^d:p(x)\leq1\},
    \qquad
    v(x)=(p(x)-1)_+.
    \label{eq:ctr-p-v}
\end{align}
Before sampling in round $t$, the learner chooses an $\F_{t-1}$-measurable
vector
\begin{align}
    s_t\in\partial v(\mu_t),
    \qquad
    s_t=0\quad\text{whenever }v(\mu_t)=0.
    \label{eq:ctr-selector}
\end{align}
The predictable measurability in \eqref{eq:ctr-selector} is an explicit oracle
assumption.  In particular, $s_t$ is fixed before $X_t$ and $\xi_t$ are
observed.

Define
\begin{align}
    \mathcal B_t(x)
     & =v(x)-v(\mu_t)-\ip{s_t}{x-\mu_t},
    \label{eq:ctr-B-definition}
    \\
    \widetilde Y_t
     & =Y_t-\frac2\eps
    \left(
    v(\mu_t)+\ip{s_t}{X_t-\mu_t}
    \right),
    \label{eq:ctr-Y-definition}
\end{align}
and
\begin{align}
    \widehat g_t^{\mathrm{ctr}}
     & =\widetilde Y_t\Sigma_t^{-1/2}Z_t+\frac2\eps s_t,
    \label{eq:ctr-g-definition}
    \\
    \widehat H_t^{\mathrm{ctr}}
     & =\widetilde Y_t\Sigma_t^{-1/2}
    (Z_tZ_t^\top-I_d)\Sigma_t^{-1/2}.
    \label{eq:ctr-H-definition}
\end{align}
The sampled quadratic is
\begin{align}
    \widehat q_t^{\mathrm{ctr}}(x)
    =\ip{\widehat g_t^{\mathrm{ctr}}}{x-\mu_t}
    +\frac\lambda2
    (x-\mu_t)^\top\widehat H_t^{\mathrm{ctr}}(x-\mu_t).
    \label{eq:ctr-qhat-definition}
\end{align}

\begin{algorithm}[!t]
    \caption{Tangent-centred unprojected Hermite--ONS}
    \label{alg:ctr-unprojected}
    \small
    \begin{tabbing}
        \qquad\=\qquad\=\kill
        \textbf{input:} $\eta,\lambda,\sigma^2,\eps$\\
        $\mu_1=0$ and $\Sigma_1=\sigma^2I_d$\\[1mm]
        \textbf{for} $t=1,\ldots,n$ \textbf{do}\\
        \>Choose predictable $s_t$ as in \eqref{eq:ctr-selector}\\
        \>Sample $Z_t\sim\N(0,I_d)$ conditionally on $\F_{t-1}$ and set
        $X_t=\mu_t+\Sigma_t^{1/2}Z_t$\\
        \>Play $A_t=X_t/\pi_\eps(X_t)$ and observe $f(A_t)+\xi_t$\\
        \>Form $Y_t$, $\mathcal B_t(X_t)$, $\widetilde Y_t$,
        $\widehat g_t^{\mathrm{ctr}}$, and $\widehat H_t^{\mathrm{ctr}}$\\
        \>\textbf{if}
        $\Sigma_t^{-1}+\lambda\eta\widehat H_t^{\mathrm{ctr}}\succ0$:\\
        \>\>Set
        $\Sigma_{t+1}^{-1}
            =\Sigma_t^{-1}+\lambda\eta\widehat H_t^{\mathrm{ctr}}$\\
        \>\>and
        $\mu_{t+1}=\mu_t-\eta\Sigma_{t+1}
            \widehat g_t^{\mathrm{ctr}}$\\
        \>\textbf{else}:\\
        \>\>Set $\Sigma_{t+1}=\Sigma_t$ and $\mu_{t+1}=\mu_t$
        \qquad\textit{$\triangleright$ The algorithm has failed.}\\
        \textbf{end for}
    \end{tabbing}
\end{algorithm}

\subsubsection{Geometry, measurability, and unbiasedness}

\begin{lemma}[Tangent geometry of the gauge violation]
    \label{lem:ctr-tangent-geometry}
    Fix $\mu\in\R^d$, let $m=v(\mu)$, and choose
    $s\in\partial v(\mu)$, taking $s=0$ when $m=0$.  Put
    \begin{align}
        \mathcal B(x)=v(x)-m-\ip{s}{x-\mu}.
    \end{align}
    Then:
    \begin{enumerate}
        \item $\mathcal B\geq0$ and $D^2\mathcal B=D^2v$ in the sense of
              distributions.
        \item For every $w\in\R^d$,
              \begin{align}
                  \ip{s}{w}              & \leq p(w),
                  \label{eq:ctr-recession-subgradient}
                  \\
                  v(\mu+w)               & \leq m+p(w),
                  \label{eq:ctr-v-increment}
                  \\
                  0\leq\mathcal B(\mu+w) & \leq p(w)+p(-w).
                  \label{eq:ctr-B-pointwise}
              \end{align}
        \item If $m>0$, then
              \begin{align}
                  p(\mu)=1+m,
                  \qquad
                  \ip{s}{\mu}=1+m,
                  \label{eq:ctr-Euler}
              \end{align}
              and for every $x\in\K_\eps$,
              \begin{align}
                  \ip{s}{\mu-x}\geq m.
                  \label{eq:ctr-inward-support}
              \end{align}
              Equality in \eqref{eq:ctr-inward-support} holds for
              $x=\mu/(1+m)$.
        \item
              \begin{align}
                  \norm{s}\leq\frac1{1-\eps}\leq2.
                  \label{eq:ctr-s-Euclidean}
              \end{align}
    \end{enumerate}
\end{lemma}

\begin{proof}
    The subgradient inequality gives
    \begin{align}
        v(x)\geq v(\mu)+\ip{s}{x-\mu},
    \end{align}
    proving $\mathcal B\geq0$.  Subtracting an affine function does not change
    a distributional Hessian, so $D^2\mathcal B=D^2v$.

    Fix $w\in\R^d$.  Subadditivity of the gauge $p$ gives, for $r>0$,
    \begin{align}
        rp(w)-p(-\mu)
        \leq p(\mu+rw)
        \leq p(\mu)+rp(w).
    \end{align}
    Hence
    \begin{align}
        \lim_{r\to\infty}\frac{v(\mu+rw)}r=p(w).
    \end{align}
    The subgradient inequality also gives
    \begin{align}
        \frac{v(\mu+rw)-v(\mu)}r\geq\ip{s}{w}.
    \end{align}
    Letting $r\to\infty$ proves
    \eqref{eq:ctr-recession-subgradient}.  Moreover,
    \begin{align}
        v(\mu+w)
         & =(p(\mu+w)-1)_+
        \leq(p(\mu)+p(w)-1)_+
        \leq m+p(w),
    \end{align}
    which is \eqref{eq:ctr-v-increment}.  Therefore,
    \begin{align}
        \mathcal B(\mu+w)
         & =v(\mu+w)-m-\ip{s}{w}
        \leq p(w)-\ip{s}{w}
        \leq p(w)+p(-w),
    \end{align}
    proving \eqref{eq:ctr-B-pointwise}.

    Suppose $m>0$.  Then $p(\mu)>1$, so $v=p-1$ in a neighbourhood of
    $\mu$ and $s\in\partial p(\mu)$.  Applying the subgradient inequality for
    $p$ with $x=0$ and $x=2\mu$ gives
    \begin{align}
        0=p(0)          & \geq p(\mu)-\ip{s}{\mu},
        \\
        2p(\mu)=p(2\mu) & \geq p(\mu)+\ip{s}{\mu}.
    \end{align}
    Thus $\ip{s}{\mu}=p(\mu)=1+m$, proving \eqref{eq:ctr-Euler}.  If
    $x\in\K_\eps$, then $v(x)=0$, and
    \begin{align}
        0=v(x)\geq m+\ip{s}{x-\mu},
    \end{align}
    which is \eqref{eq:ctr-inward-support}.  For
    $x=\mu/(1+m)$, \eqref{eq:ctr-Euler} gives $\ip{s}{x}=1$, so equality
    holds.  Finally, a subgradient of a finite $L$-Lipschitz convex function
    has Euclidean norm at most $L$.  Since $v$ is
    $(1-\eps)^{-1}$-Lipschitz, \eqref{eq:ctr-s-Euclidean} follows.
\end{proof}

\begin{lemma}[Conditional measurability, integrability, and unbiasedness]
    \label{lem:ctr-unbiasedness}
    On every active round, $\mathcal B_t(X_t)$, $\widetilde Y_t$,
    $\widehat g_t^{\mathrm{ctr}}$, and $\widehat H_t^{\mathrm{ctr}}$ are
    $\F_t$-measurable.  Every entry has a finite conditional moment of every
    finite order.  Moreover,
    \begin{align}
        \E_t[\widehat g_t^{\mathrm{ctr}}] & =g_t,
                                          &
        \E_t[\widehat H_t^{\mathrm{ctr}}] & =H_t,
        \label{eq:ctr-unbiasedness}
    \end{align}
    and, for every $\F_{t-1}$-measurable $x\in\R^d$,
    \begin{align}
        \E_t[\widehat q_t^{\mathrm{ctr}}(x)]=q_t(x).
        \label{eq:ctr-q-unbiasedness}
    \end{align}
\end{lemma}

\begin{proof}
    The maps $\pi$, $p$, and $v$ are continuous.  The quantities
    $\mu_t$, $\Sigma_t$, and $s_t$ are $\F_{t-1}$-measurable, while
    $X_t$, $Z_t$, and the observed feedback are $\F_t$-measurable.  Thus all
    quantities in \eqref{eq:ctr-B-definition}--\eqref{eq:ctr-H-definition}
    are $\F_t$-measurable.

    Conditionally on $\F_{t-1}$, write $W_t=X_t-\mu_t$.  By
    \eqref{eq:ctr-v-increment} and \eqref{eq:ctr-B-pointwise},
    \begin{align}
        \pi_\eps(X_t)     & \leq1+v(\mu_t)+p(W_t),
        \\
        \mathcal B_t(X_t) & \leq p(W_t)+p(-W_t).
        \label{eq:ctr-integrability-envelope}
    \end{align}
    The gauge $p$ is Lipschitz, so the right-hand sides have finite Gaussian
    moments of every finite order.  The noise assumption gives, for every
    $r>0$,
    \begin{align}
        \E_t[|\xi_t|^r\mid X_t]<\infty,
    \end{align}
    because $|z|^r\leq C_r e^{z^2}$ for a finite deterministic $C_r$.
    Equation \eqref{eq:ctr-Y-definition}, Holder's inequality, and Gaussian
    moments therefore imply the asserted conditional integrability of the
    two estimators.

    Define the predictable affine function
    \begin{align}
        a_t(x)=\frac2\eps
        \left(v(\mu_t)+\ip{s_t}{x-\mu_t}\right).
    \end{align}
    Conditional centring of the noise gives
    \begin{align}
        \E_t[\widetilde Y_t\mid X_t]=e(X_t)-a_t(X_t).
        \label{eq:ctr-Y-conditional-mean}
    \end{align}
    Since $a_t$ is affine,
    \begin{align}
        D(e-a_t)   & =De-\frac2\eps s_t,
                   &
        D^2(e-a_t) & =D^2e.
    \end{align}
    Distributional Gaussian integration by parts, justified by the preceding
    integrability bounds, yields
    \begin{align}
        \E_t\left[
            \widetilde Y_t\Sigma_t^{-1/2}Z_t
            \right]
         & =\left\langle D(e-a_t),\gamma_{\mu_t,\Sigma_t}\right\rangle
        =g_t-\frac2\eps s_t,
        \\
        \E_t\left[
            \widetilde Y_t\Sigma_t^{-1/2}
            (Z_tZ_t^\top-I_d)\Sigma_t^{-1/2}
            \right]
         & =\left\langle D^2(e-a_t),\gamma_{\mu_t,\Sigma_t}\right\rangle
        =H_t.
    \end{align}
    Adding the deterministic correction in
    \eqref{eq:ctr-g-definition} proves \eqref{eq:ctr-unbiasedness};
    \eqref{eq:ctr-q-unbiasedness} follows by linearity.
\end{proof}

\begin{lemma}[Tangent-remainder moments and normal width]
    \label{lem:ctr-B-moments}
    Fix a round and condition on $\F_{t-1}$.  Suppose
    \begin{align}
        \Sigma_t\preceq2\sigma^2I_d,
        \qquad
        \sigma=\frac1{M\sqrt{2d}},
        \qquad
        0<\eps\leq\frac12.
        \label{eq:ctr-B-cov-assumption}
    \end{align}
    Put
    \begin{align}
        V_t=\left\langle D^2v,\gamma_{\mu_t,\Sigma_t}\right\rangle,
        \qquad
        c_t=\tr(\Sigma_tV_t),
        \qquad
        m_t=v(\mu_t).
    \end{align}
    Then
    \begin{align}
        0\leq c_t\leq2,
        \qquad
        \E_t[\mathcal B_t(X_t)]\leq c_t
        \leq\eps\tr(H_t\Sigma_t).
        \label{eq:ctr-B-first}
    \end{align}
    For every $u\geq0$,
    \begin{align}
        \Pp_t\left(\mathcal B_t(X_t)>4+2u\right)
        \leq\exp\left(-\frac{u^2}{8}\right).
        \label{eq:ctr-B-tail}
    \end{align}
    Hence, for every $\ell\geq1$,
    \begin{align}
        \E_t[\mathcal B_t(X_t)^2]
         & \leq32\sqrt\ell\,c_t+e^{-8\ell}
        \label{eq:ctr-B-second-a}
        \\
         & \leq32\eps\sqrt\ell\,\tr(H_t\Sigma_t)+e^{-8\ell}.
        \label{eq:ctr-B-second-b}
    \end{align}
    Finally,
    \begin{align}
        s_t^\top\Sigma_ts_t
        \leq4m_t^2+2^{11}c_t
        \leq4m_t^2+2^{11}\eps\tr(H_t\Sigma_t).
        \label{eq:ctr-normal-width}
    \end{align}
\end{lemma}

\begin{proof}
    Write $W_t=X_t-\mu_t$.  At every differentiability point of $v$, the
    supporting-hyperplane inequality at $X_t$ gives
    \begin{align}
        v(X_t)-v(\mu_t)
        \leq\ip{\nabla v(X_t)}{W_t}.
    \end{align}
    Since $\E_t[W_t]=0$, boundedness of the almost-everywhere gradient and
    distributional integration by parts yield
    \begin{align}
        \E_t[\mathcal B_t(X_t)]
         & =\E_t[v(X_t)]-v(\mu_t)
        \\
         & \leq\E_t\left[\ip{\nabla v(X_t)}{W_t}\right]
        =\tr(\Sigma_tV_t)=c_t.
        \label{eq:ctr-B-IBP}
    \end{align}
    The extension decomposition
    \begin{align}
        e=\bar e+\frac{1+\eps}{\eps}v
    \end{align}
    and convexity of $\bar e$ imply
    \begin{align}
        D^2v\preceq\frac\eps{1+\eps}D^2e\preceq\eps D^2e.
    \end{align}
    Pairing with the Gaussian density and taking the trace after multiplication
    by $\Sigma_t$ proves the last inequality in \eqref{eq:ctr-B-first}.

    Let $r(x)$ denote an almost-everywhere gradient of $v$.  Applying
    \eqref{eq:ctr-recession-subgradient} at $X_t$ gives
    \begin{align}
        \ip{r(X_t)}{W_t}\leq p(W_t).
    \end{align}
    Under \eqref{eq:ctr-B-cov-assumption}, take an independent
    $Q_t\sim\N(0,2\sigma^2I_d-\Sigma_t)$.  Jensen's inequality and the
    definition of $M$ give
    \begin{align}
        \E_t[\pi(W_t)]
         & =\E_t\left[\pi\left(\E_t[W_t+Q_t\mid W_t]\right)\right]
        \\
         & \leq\E_t[\pi(W_t+Q_t)]
        =\sqrt2\sigma\,\E[\norm Z]\E[\pi(U)]
        \leq\sqrt{2d}\,\sigma M=1.
        \label{eq:ctr-pi-mean}
    \end{align}
    Therefore
    \begin{align}
        c_t
        =\E_t\left[\ip{r(X_t)}{W_t}\right]
        \leq\E_t[p(W_t)]\leq2.
    \end{align}

    Set $S_t=\pi(W_t)+\pi(-W_t)$.  Symmetry gives
    $\E_t[S_t]\leq2$, and $S_t$ is $2$-Lipschitz.  Since
    $\norm{\Sigma_t}_{\op}\leq2\sigma^2\leq1$, Gaussian concentration gives
    \begin{align}
        \Pp_t(S_t>2+u)\leq e^{-u^2/8}.
        \label{eq:ctr-S-tail}
    \end{align}
    Equation \eqref{eq:ctr-B-pointwise} and
    $p\leq2\pi$ imply $\mathcal B_t(X_t)\leq2S_t$, so
    \eqref{eq:ctr-B-tail} follows.

    Put $A=32\sqrt\ell$.  Then
    \begin{align}
        \E_t[\mathcal B_t(X_t)^2]
         & \leq
        A\E_t[\mathcal B_t(X_t)]
        +\E_t\left[
            \mathcal B_t(X_t)^2\1\{\mathcal B_t(X_t)>A\}
            \right].
        \label{eq:ctr-B-split}
    \end{align}
    For $z\geq A$, \eqref{eq:ctr-B-tail} gives
    \begin{align}
        \Pp_t(\mathcal B_t(X_t)>z)
        \leq\exp\left(-\frac{(z-4)^2}{32}\right).
    \end{align}
    The layer-cake identity and
    $\int_a^\infty e^{-u^2/32}\,du\leq16a^{-1}e^{-a^2/32}$ give
    \begin{align}
         & \E_t\left[
            \mathcal B_t(X_t)^2\1\{\mathcal B_t(X_t)>A\}
            \right]
        \\
         & \quad=
        A^2\Pp_t(\mathcal B_t(X_t)>A)
        +\int_A^\infty2z\Pp_t(\mathcal B_t(X_t)>z)\,dz
        \\
         & \quad\leq
        \left(A^2+32+\frac{128}{A-4}\right)
        e^{-(A-4)^2/32}
        \leq e^{-8\ell}.
    \end{align}
    Combining this with \eqref{eq:ctr-B-split} and
    \eqref{eq:ctr-B-first} proves
    \eqref{eq:ctr-B-second-a}--\eqref{eq:ctr-B-second-b}.

    It remains to prove \eqref{eq:ctr-normal-width}.  Let
    \begin{align}
        T_t=\ip{s_t}{W_t},
        \qquad
        a_t=s_t^\top\Sigma_ts_t.
    \end{align}
    Conditionally, $T_t\sim\N(0,a_t)$.  Since $v(X_t)\geq0$,
    \begin{align}
        \mathcal B_t(X_t)
        =v(X_t)-m_t-T_t
        \geq(-m_t-T_t)_+.
        \label{eq:ctr-shortfall}
    \end{align}
    If $m_t\geq\sqrt{a_t}/2$, then $a_t\leq4m_t^2$.  Otherwise, on
    $\{T_t\leq-\sqrt{a_t}\}$,
    \begin{align}
        (-m_t-T_t)_+\geq\frac12\sqrt{a_t}.
    \end{align}
    Moreover,
    \begin{align}
        \Pp_t(T_t\leq-\sqrt{a_t})
        =\Pp(Z\leq-1)
        \geq\int_1^{3/2}\frac{e^{-u^2/2}}{\sqrt{2\pi}}\,du
        >\frac1{16}.
    \end{align}
    Thus, in the second case,
    \begin{align}
        c_t\geq\E_t[\mathcal B_t(X_t)]\geq\frac1{32}\sqrt{a_t},
    \end{align}
    so $a_t\leq2^{10}c_t^2\leq2^{11}c_t$, because $c_t\leq2$.
    Combining the two cases proves \eqref{eq:ctr-normal-width}.
\end{proof}

\subsubsection{Stopped invariants and coarse confinement}

Retain the predictable covariance
\begin{align}
    \overline\Sigma_1^{-1}=\Sigma_1^{-1}=\sigma^{-2}I_d,
    \qquad
    \overline\Sigma_{t+1}^{-1}
    =\overline\Sigma_t^{-1}+\lambda\eta H_t.
    \label{eq:ctr-predictable-covariance}
\end{align}
Define $\tau_{\mathrm{cov}}$ by replacing $\widehat H_t$ with
$\widehat H_t^{\mathrm{ctr}}$ in the covariance stopping time of the
manuscript.  Thus, for $t\leq\tau_{\mathrm{cov}}$,
\begin{align}
    \frac12\overline\Sigma_t^{-1}
    \preceq\Sigma_t^{-1}
    \preceq\frac32\overline\Sigma_t^{-1}.
    \label{eq:ctr-cov-comparison}
\end{align}
Fix
\begin{align}
    B=\frac{b}{\lambda^2\ell}
    \label{eq:ctr-focus-radius}
\end{align}
and define
\begin{align}
    \tau_{\mathrm{focus}}
    =\min\left(
    \left\{
    t\in[n]:
    \frac12
    \normm{\mu_{t+1}-x_\star}{\Sigma_{t+1}^{-1}}^2>B
    \right\}\cup\{n\}
    \right),
    \qquad
    \tau=\tau_{\mathrm{cov}}\wedge\tau_{\mathrm{focus}}.
    \label{eq:ctr-safe-stop}
\end{align}
Finally, set
\begin{align}
    \rho=4\sigma\sqrt B
    =\frac{2\sqrt{2b}}{c_\lambda}\,d\ell^{3/2}
    \label{eq:ctr-rho}
\end{align}
when $\lambda=c_\lambda/(d^{3/2}M\ell^2)$.

\begin{lemma}[Stopped distance and coarse confinement]
    \label{lem:ctr-coarse-confinement}
    Suppose $t\leq\tau$.  Define the distance-to-$\K_\eps$ potential
    \begin{align}
        \mathsf D_t
        =\frac12\inf_{x\in\K_\eps}
        \normm{\mu_t-x}{\Sigma_t^{-1}}^2.
        \label{eq:ctr-distance-potential}
    \end{align}
    Then
    \begin{align}
        \mathsf D_t
        \leq\frac12
        \normm{\mu_t-x_\star}{\Sigma_t^{-1}}^2
        \leq B,
        \label{eq:ctr-distance-below-focus}
    \end{align}
    and
    \begin{align}
        v(\mu_t)\leq4\sigma\sqrt{\mathsf D_t}
        \leq\rho.
        \label{eq:ctr-coarse-v-bound}
    \end{align}
    Moreover,
    \begin{align}
        \Sigma_t\preceq2\overline\Sigma_t
        \preceq2\sigma^2I_d.
        \label{eq:ctr-stopped-cov-upper}
    \end{align}
\end{lemma}

\begin{proof}
    Since $x_\star\in\K_\eps$, the first inequality in
    \eqref{eq:ctr-distance-below-focus} follows from the definition of
    $\mathsf D_t$, and the second follows from the focus stop.

    Let $P_t$ be the unique minimizer in
    \eqref{eq:ctr-distance-potential}.  Since $p(P_t)\leq1$, subadditivity of
    $p$ gives
    \begin{align}
        v(\mu_t)
         & =(p(\mu_t)-1)_+
        \leq p(\mu_t-P_t)
        \leq2\pi(\mu_t-P_t)
        \leq2\norm{\mu_t-P_t}.
        \label{eq:ctr-v-to-Euclidean-distance}
    \end{align}
    The lower covariance comparison and monotonicity of
    $\overline\Sigma_t^{-1}$ give
    \begin{align}
        \Sigma_t^{-1}
        \succeq\frac12\overline\Sigma_t^{-1}
        \succeq\frac1{2\sigma^2}I_d.
    \end{align}
    Hence
    \begin{align}
        \norm{\mu_t-P_t}^2
        \leq2\sigma^2
        \normm{\mu_t-P_t}{\Sigma_t^{-1}}^2
        =4\sigma^2\mathsf D_t.
    \end{align}
    Combining this with \eqref{eq:ctr-v-to-Euclidean-distance} proves
    \eqref{eq:ctr-coarse-v-bound}.  Inverting the lower covariance comparison
    and using $\overline\Sigma_t\preceq\overline\Sigma_1=\sigma^2I_d$ proves
    \eqref{eq:ctr-stopped-cov-upper}.
\end{proof}


\begin{lemma}[Predictable spectrum and trace budget]
    \label{lem:ctr-trace}
    Let $L=\Lip(e)$ and suppose
    \begin{align}
        \lambda\eta L\leq1,
        \qquad
        \sqrt{\frac32}\lambda\eta L\sigma\leq1.
        \label{eq:ctr-trace-hypotheses}
    \end{align}
    Then, for every $t\leq\tau_{\mathrm{cov}}$,
    \begin{align}
        \norm{\overline\Sigma_{t+1}^{-1}}_{\op}^{1/2}
        \leq\sigma^{-1}+t\lambda\eta L
        \leq\sigma^{-1}+t,
        \label{eq:ctr-trace-spectrum}
    \end{align}
    and
    \begin{align}
        \eta\sum_{t=1}^{\tau_{\mathrm{cov}}}
        \tr(H_t\overline\Sigma_t)
        \leq\frac{4d\log(1+n)}\lambda.
        \label{eq:ctr-trace-logdet}
    \end{align}
\end{lemma}

\begin{proof}
    Distributional integration by parts gives, for every Euclidean unit vector
    $u$,
    \begin{align}
        u^\top H_tu
         & =\E_t\left[
            \ip{\nabla e(X_t)}{u}
            \ip{\Sigma_t^{-1}(X_t-\mu_t)}{u}
            \right]
        \\
         & \leq L\sqrt{u^\top\Sigma_t^{-1}u}.
    \end{align}
    Hence
    \begin{align}
        \norm{H_t}_{\op}
        \leq L\norm{\Sigma_t^{-1}}_{\op}^{1/2}
        \leq\sqrt{\frac32}L
        \norm{\overline\Sigma_t^{-1}}_{\op}^{1/2},
        \label{eq:ctr-H-op-bound}
    \end{align}
    where covariance comparison is used in the last step.  Therefore
    \begin{align}
        \norm{\overline\Sigma_{t+1}^{-1}}_{\op}
         & \leq
        \norm{\overline\Sigma_t^{-1}}_{\op}
        +\sqrt{\frac32}\lambda\eta L
        \norm{\overline\Sigma_t^{-1}}_{\op}^{1/2}
        \\
         & \leq
        \left(
        \norm{\overline\Sigma_t^{-1}}_{\op}^{1/2}
        +\lambda\eta L
        \right)^2.
    \end{align}
    Taking square roots and iterating proves
    \eqref{eq:ctr-trace-spectrum}.

    We next bound the normalised predictable Hessian.  For a Euclidean unit
    vector $u$, the first inequality in the proof, applied to
    $\overline\Sigma_t^{1/2}u$, gives
    \begin{align}
        u^\top\overline\Sigma_t^{1/2}H_t
        \overline\Sigma_t^{1/2}u
         & \leq
        L\norm{\overline\Sigma_t^{1/2}u}
        \sqrt{
            u^\top\overline\Sigma_t^{1/2}
            \Sigma_t^{-1}
            \overline\Sigma_t^{1/2}u
        }
        \\
         & \leq\sqrt{\frac32}L
        \norm{\overline\Sigma_t^{1/2}u}
        \leq\sqrt{\frac32}L\sigma.
    \end{align}
    Thus every eigenvalue of
    \begin{align}
        C_t=\lambda\eta
        \overline\Sigma_t^{1/2}H_t\overline\Sigma_t^{1/2}
    \end{align}
    belongs to $[0,1]$.  The scalar inequality
    $x\leq2\log(1+x)$ for $x\in[0,1]$ and cyclicity of the trace give
    \begin{align}
        \eta\sum_{t=1}^{\tau_{\mathrm{cov}}}
        \tr(H_t\overline\Sigma_t)
         & =\frac1\lambda
        \sum_{t=1}^{\tau_{\mathrm{cov}}}\tr(C_t)
        \\
         & \leq\frac2\lambda
        \sum_{t=1}^{\tau_{\mathrm{cov}}}
        \log\det(I_d+C_t)
        \\
         & =\frac2\lambda
        \log\frac{
            \det(\overline\Sigma_{\tau_{\mathrm{cov}}+1}^{-1})
        }{
            \det(\overline\Sigma_1^{-1})
        }
        \\
         & \leq\frac{2d}\lambda
        \log\left(
        \sigma^2
        \norm{\overline\Sigma_{\tau_{\mathrm{cov}}+1}^{-1}}_{\op}
        \right)
        \\
         & =\frac{4d}\lambda
        \log\left(
        \sigma
        \norm{\overline\Sigma_{\tau_{\mathrm{cov}}+1}^{-1}}_{\op}^{1/2}
        \right)
        \\
         & \leq\frac{4d\log(1+n)}\lambda,
    \end{align}
    where the last inequality uses
    \eqref{eq:ctr-trace-spectrum}, $\sigma\leq1$, and
    $\tau_{\mathrm{cov}}\leq n$.
\end{proof}

\subsubsection{Centred feedback and weighted Gaussian scores}

Put
\begin{align}
    R_\rho=1+\rho,
    \qquad
    B_Y=2^8\ell\left(R_\rho+\eps^{-1}\right).
    \label{eq:ctr-feedback-scales}
\end{align}

\begin{lemma}[Predictable and realised centred feedback]
    \label{lem:ctr-feedback}
    Suppose $\ell\geq64$ and work through the stopping time $\tau$ in
    \eqref{eq:ctr-safe-stop}.  Suppose also that
    \begin{align}
        \sum_{t=1}^{\tau}\tr(H_t\Sigma_t)
        \leq\frac{8d\log(1+n)}{\lambda\eta},
        \label{eq:ctr-feedback-trace-input}
    \end{align}
    and that
    \begin{align}
        \frac{2^{12}d\sqrt\ell\log(1+n)}
        {\lambda\eta\eps}
             & \leq nR_\rho^2,
        \label{eq:ctr-feedback-param-a}
        \\
        \frac{ne^{-8\ell}}{\eps^2}
             & \leq nR_\rho^2,
        \label{eq:ctr-feedback-param-b}
        \\
        B_Y\sqrt\ell
             & \leq2^{-8}\sqrt n\,R_\rho,
        \label{eq:ctr-feedback-param-c}
        \\
        \ell & \geq\log\frac{512n}{\delta}.
        \label{eq:ctr-feedback-param-d}
    \end{align}
    Then
    \begin{align}
        \sum_{t=1}^{\tau}\E_t[\widetilde Y_t^2]
        \leq2^{12}nR_\rho^2.
        \label{eq:ctr-predictable-Y2}
    \end{align}
    Moreover, with probability at least $1-\delta/16$,
    \begin{align}
        \max_{t\leq\tau}|\widetilde Y_t|
         & \leq B_Y,
        \label{eq:ctr-Y-maximum}
        \\
        \sum_{t=1}^{\tau}\widetilde Y_t^2
         & \leq2^{13}nR_\rho^2.
        \label{eq:ctr-realised-Y2}
    \end{align}
\end{lemma}

\begin{proof}
    Write $r_t=\pi_\eps(X_t)$ and $W_t=X_t-\mu_t$.  From
    \eqref{eq:ctr-Y-definition},
    \begin{align}
        \widetilde Y_t
        =r_t\bigl(f(A_t)+\xi_t\bigr)
        +\frac2\eps\mathcal B_t(X_t).
        \label{eq:ctr-Y-decomposition}
    \end{align}
    Conditioning first on $X_t$, using
    $\E_t[\xi_t\mid X_t]=0$ and
    $\E_t[\xi_t^2\mid X_t]\leq1$, gives
    \begin{align}
        \E_t[\widetilde Y_t^2\mid X_t]
         & \leq3r_t^2+\frac8{\eps^2}\mathcal B_t(X_t)^2.
        \label{eq:ctr-Y-conditional-second}
    \end{align}
    By \cref{lem:ctr-coarse-confinement} and
    \eqref{eq:ctr-v-increment},
    \begin{align}
        r_t\leq R_\rho+p(W_t).
        \label{eq:ctr-r-envelope}
    \end{align}
    Let $\Delta_t=\pi(W_t)-\E_t[\pi(W_t)]$.  Gaussian concentration and
    \eqref{eq:ctr-stopped-cov-upper} give
    \begin{align}
        \Pp_t(|\Delta_t|>u)
        \leq2e^{-u^2/2},
        \qquad u\geq0,
    \end{align}
    because $\norm{\Sigma_t}_{\op}\leq1$.  Therefore
    \begin{align}
        \E_t[\Delta_t^2]
         & =\int_0^\infty2u\Pp_t(|\Delta_t|>u)\,du
        \leq4.
    \end{align}
    Together with \eqref{eq:ctr-pi-mean}, this yields
    \begin{align}
        \E_t[\pi(W_t)^2]
        \leq2\E_t[\pi(W_t)]^2+2\E_t[\Delta_t^2]
        \leq10.
    \end{align}
    Hence, since $p\leq2\pi$,
    \begin{align}
        \E_t[r_t^2]
        \leq2R_\rho^2+2\E_t[p(W_t)^2]
        \leq82R_\rho^2
        \leq2^7R_\rho^2.
        \label{eq:ctr-r-second}
    \end{align}
    Applying \cref{lem:ctr-B-moments} in
    \eqref{eq:ctr-Y-conditional-second} gives
    \begin{align}
        \E_t[\widetilde Y_t^2]
        \leq
        2^9R_\rho^2
        +\frac{2^8\sqrt\ell}{\eps}\tr(H_t\Sigma_t)
        +\frac{8e^{-8\ell}}{\eps^2}.
        \label{eq:ctr-Y-roundwise-second}
    \end{align}
    Summing, using \eqref{eq:ctr-feedback-trace-input}, and then applying
    \eqref{eq:ctr-feedback-param-a}--\eqref{eq:ctr-feedback-param-b} proves
    \eqref{eq:ctr-predictable-Y2}.

    We next prove the pointwise truncation bound.  Gaussian concentration,
    \eqref{eq:ctr-pi-mean}, symmetry of $W_t$, and the noise assumption give
    \begin{align}
        \Pp_t\bigl(\pi(W_t)>1+2\sqrt\ell\bigr)  & \leq e^{-2\ell},
        \\
        \Pp_t\bigl(\pi(-W_t)>1+2\sqrt\ell\bigr) & \leq e^{-2\ell},
        \\
        \Pp_t(|\xi_t|>\sqrt\ell\mid X_t)        & \leq2e^{-\ell}.
        \label{eq:ctr-roundwise-tail-events}
    \end{align}
    Outside these three events,
    \begin{align}
        r_t
         & \leq R_\rho+2+4\sqrt\ell
        \leq8R_\rho\sqrt\ell,
        \\
        \mathcal B_t(X_t)
         & \leq p(W_t)+p(-W_t)
        \leq4+8\sqrt\ell
        \leq12\sqrt\ell.
    \end{align}
    Thus
    \begin{align}
        |\widetilde Y_t|
         & \leq r_t(1+|\xi_t|)+\frac2\eps\mathcal B_t(X_t)
        \\
         & \leq16R_\rho\ell+\frac{24\sqrt\ell}{\eps}
        \leq B_Y.
        \label{eq:ctr-Y-pointwise}
    \end{align}
    Since $\1\{t\leq\tau\}$ is $\F_{t-1}$-measurable, a union bound and
    \eqref{eq:ctr-feedback-param-d} show that
    \eqref{eq:ctr-Y-maximum} fails with probability at most $\delta/32$.

    Define
    \begin{align}
        \overline Y_t
        =\widetilde Y_t\1\{|\widetilde Y_t|\leq B_Y\},
        \qquad
        D_t=\1\{t\leq\tau\}
        \left(
        \overline Y_t^2-\E_t[\overline Y_t^2]
        \right).
    \end{align}
    Then $\E_t[D_t]=0$, $|D_t|\leq B_Y^2$, and
    \begin{align}
        \sum_{t=1}^n\E_t[D_t^2]
         & \leq B_Y^2
        \sum_{t=1}^{\tau}\E_t[\widetilde Y_t^2]
        \leq2^{12}nB_Y^2R_\rho^2.
    \end{align}
    Apply the manuscript's Freedman inequality with
    \begin{align}
        b=B_Y^2,
        \qquad
        s=2^6B_Y\sqrt n\,R_\rho,
        \qquad
        z=\sqrt{2\log\frac{64}{\delta}}\leq\sqrt{2\ell}.
    \end{align}
    Outside an event of probability $\delta/32$,
    \begin{align}
        \sum_{t=1}^nD_t
         & \leq sz+\frac{bz^2}{3}
        \\
         & \leq2^7B_YR_\rho\sqrt{n\ell}+B_Y^2\ell
        \leq2^{12}nR_\rho^2,
    \end{align}
    where the last inequality is
    \eqref{eq:ctr-feedback-param-c}.  On the event
    \eqref{eq:ctr-Y-maximum}, $\overline Y_t=\widetilde Y_t$ for every
    $t\leq\tau$.  Combining the last display with
    \eqref{eq:ctr-predictable-Y2} proves
    \eqref{eq:ctr-realised-Y2}.  A union bound completes the proof.
\end{proof}

\begin{lemma}[Weighted centred Gaussian scores]
    \label{lem:ctr-score-localisation}
    Let $t\leq\tau$, let $h_t$ be an $\F_{t-1}$-measurable unit vector, and
    put $G_t=\ip{h_t}{Z_t}$.  Suppose
    \begin{align}
        2^{20}\eps^{-2}R_\rho^2\leq e^{4\ell}.
        \label{eq:ctr-score-tail-parameter}
    \end{align}
    Then
    \begin{align}
        \E_t\left[
            \widetilde Y_t^2(G_t^2+1)^2
            \right]
        \leq2^{14}\ell^2\E_t[\widetilde Y_t^2]+e^{-4\ell}.
        \label{eq:ctr-score-localisation}
    \end{align}
\end{lemma}

\begin{proof}
    Let $\mathcal A_t=\{|G_t|\leq8\sqrt\ell\}$.  On $\mathcal A_t$,
    $(G_t^2+1)^2\leq2^{14}\ell^2$.  It remains to control the complement.

    Set
    \begin{align}
        \mathcal R_t
        =R_\rho+p(W_t)+p(-W_t).
    \end{align}
    Equations \eqref{eq:ctr-r-envelope}, \eqref{eq:ctr-B-pointwise}, and
    \eqref{eq:ctr-Y-conditional-second} give
    \begin{align}
        \E_t[\widetilde Y_t^2\mid X_t]
        \leq11\eps^{-2}\mathcal R_t^2.
        \label{eq:ctr-score-envelope}
    \end{align}
    Put $S_t=\pi(W_t)+\pi(-W_t)$.  As in
    \eqref{eq:ctr-S-tail},
    \begin{align}
        \Pp_t(|S_t-\E_tS_t|>u)\leq2e^{-u^2/8}.
    \end{align}
    The layer-cake identity gives
    \begin{align}
        \E_t[|S_t-\E_tS_t|^4]
        \leq8\int_0^\infty u^3e^{-u^2/8}\,du
        =2^8.
    \end{align}
    Since $\E_tS_t\leq2$, Minkowski's inequality yields
    $\norm{S_t}_{L^4(\Pp_t)}\leq6$.  Because
    $p(W_t)+p(-W_t)\leq2S_t$,
    \begin{align}
        \E_t[\mathcal R_t^4]^{1/2}
        \leq2^8R_\rho^2.
        \label{eq:ctr-R-fourth}
    \end{align}
    Conditioning on $X_t$ and applying Cauchy--Schwarz gives
    \begin{align}
         & \E_t\left[
            \widetilde Y_t^2(G_t^2+1)^2\1_{\mathcal A_t^c}
            \right]
        \\
         & \quad\leq
        11\eps^{-2}\E_t[\mathcal R_t^4]^{1/2}
        \E_t\left[
            (G_t^2+1)^4\1_{\mathcal A_t^c}
            \right]^{1/2}.
        \label{eq:ctr-score-tail-CS}
    \end{align}
    Since $G_t\sim\N(0,1)$ conditionally, another Cauchy--Schwarz
    inequality, $\E[G_t^{16}]=15!!$, and
    $\Pp(|G_t|>8\sqrt\ell)\leq2e^{-32\ell}$ imply
    \begin{align}
        \E_t\left[
            (G_t^2+1)^4\1_{\mathcal A_t^c}
            \right]^{1/2}
        \leq2^8e^{-8\ell}.
    \end{align}
    Substitution into \eqref{eq:ctr-score-tail-CS} gives
    \begin{align}
        \E_t\left[
            \widetilde Y_t^2(G_t^2+1)^2\1_{\mathcal A_t^c}
            \right]
        \leq2^{20}\eps^{-2}R_\rho^2e^{-8\ell}
        \leq e^{-4\ell},
    \end{align}
    where the last inequality is
    \eqref{eq:ctr-score-tail-parameter}.  This proves
    \eqref{eq:ctr-score-localisation}.
\end{proof}

\subsubsection{Unprojected ONS and the deterministic tangent correction}

\begin{lemma}[Exact unprojected ONS identity]
    \label{lem:ctr-unprojected-ONS}
    On every round on which the update in \cref{alg:ctr-unprojected} is
    performed, for every $x\in\R^d$,
    \begin{align}
         & \frac12\left(
        \normm{\mu_{t+1}-x}{\Sigma_{t+1}^{-1}}^2
        -\normm{\mu_t-x}{\Sigma_t^{-1}}^2
        \right)
        \nonumber        \\
         & \qquad=
        \eta\widehat q_t^{\mathrm{ctr}}(x)
        +\frac{\eta^2}{2}
        \normm{\widehat g_t^{\mathrm{ctr}}}{\Sigma_{t+1}}^2.
        \label{eq:ctr-unprojected-ONS}
    \end{align}
\end{lemma}

\begin{proof}
    Put $A_{t+1}=\Sigma_{t+1}^{-1}$ and
    $\Delta_t=-\eta\Sigma_{t+1}\widehat g_t^{\mathrm{ctr}}$.  Then
    \begin{align}
        \frac12\normm{\mu_t+\Delta_t-x}{A_{t+1}}^2
         & =\frac12\normm{\mu_t-x}{A_{t+1}}^2
        +\eta\ip{\widehat g_t^{\mathrm{ctr}}}{x-\mu_t}
        +\frac{\eta^2}{2}
        \normm{\widehat g_t^{\mathrm{ctr}}}{\Sigma_{t+1}}^2
        \\
         & =\frac12\normm{\mu_t-x}{\Sigma_t^{-1}}^2
        +\frac{\lambda\eta}{2}
        (x-\mu_t)^\top\widehat H_t^{\mathrm{ctr}}(x-\mu_t)
        \\
         & \qquad
        +\eta\ip{\widehat g_t^{\mathrm{ctr}}}{x-\mu_t}
        +\frac{\eta^2}{2}
        \normm{\widehat g_t^{\mathrm{ctr}}}{\Sigma_{t+1}}^2.
    \end{align}
    Since $\mu_t+\Delta_t=\mu_{t+1}$, the right-hand side is exactly the
    right-hand side of \eqref{eq:ctr-unprojected-ONS} plus the old potential.
\end{proof}

\begin{lemma}[Radial coercivity]
    \label{lem:ctr-radial-coercivity}
    Put
    \begin{align}
        \mathfrak R_t=\E_t[e(X_t)]-e(x_\star).
    \end{align}
    Then
    \begin{align}
        \mathfrak R_t\geq\frac2\eps v(\mu_t).
        \label{eq:ctr-radial-coercivity}
    \end{align}
\end{lemma}

\begin{proof}
    Jensen's inequality gives $\E_t[e(X_t)]\geq e(\mu_t)$.  If
    $v(\mu_t)=0$, then $e(\mu_t)\geq e(x_\star)$ because $x_\star$ is a
    global minimizer of $e$.  Suppose $m_t=v(\mu_t)>0$ and put
    \begin{align}
        a_t=\frac{\mu_t}{1+m_t}\in\K_\eps.
    \end{align}
    The definition of the extension and optimality of $x_\star$ on
    $\K_\eps$ give
    \begin{align}
        e(\mu_t)-e(x_\star)
         & =(1+m_t)f(a_t)+\frac2\eps m_t-f(x_\star)
        \\
         & =f(a_t)-f(x_\star)+m_tf(a_t)+\frac2\eps m_t
        \geq\frac2\eps m_t.
    \end{align}
    Combining the two cases proves \eqref{eq:ctr-radial-coercivity}.
\end{proof}

\begin{lemma}[Centred gradient energy]
    \label{lem:ctr-gradient-energy}
    Suppose that $\Sigma_{t+1}\preceq3\Sigma_t$ for every $t\leq\tau$.
    Then
    \begin{align}
         & \sum_{t=1}^{\tau}
        \normm{\widehat g_t^{\mathrm{ctr}}}{\Sigma_{t+1}}^2
        \nonumber            \\
         & \quad\leq
        6\sum_{t=1}^{\tau}\widetilde Y_t^2\norm{Z_t}^2
        +\frac{96}{\eps^2}\sum_{t=1}^{\tau}v(\mu_t)^2
        +\frac{3\cdot2^{14}}\eps
        \sum_{t=1}^{\tau}\tr(H_t\Sigma_t).
        \label{eq:ctr-full-gradient-energy}
    \end{align}
    Consequently, on the stopped interval where $v(\mu_t)\leq\rho$,
    \begin{align}
         & \frac{\eta^2}{2}
        \sum_{t=1}^{\tau}
        \normm{\widehat g_t^{\mathrm{ctr}}}{\Sigma_{t+1}}^2
        \nonumber           \\
         & \quad\leq
        3\eta^2\sum_{t=1}^{\tau}
        \widetilde Y_t^2\norm{Z_t}^2
        +\frac{24\eta^2\rho}{\eps}
        \sum_{t=1}^{\tau}\mathfrak R_t
        +\frac{3\cdot2^{13}\eta^2}\eps
        \sum_{t=1}^{\tau}\tr(H_t\Sigma_t),
        \label{eq:ctr-focus-energy}
    \end{align}
    and
    \begin{align}
         & \frac\eta2
        \sum_{t=1}^{\tau}
        \normm{\widehat g_t^{\mathrm{ctr}}}{\Sigma_{t+1}}^2
        \nonumber     \\
         & \quad\leq
        3\eta\sum_{t=1}^{\tau}
        \widetilde Y_t^2\norm{Z_t}^2
        +\frac{24\eta\rho}{\eps}
        \sum_{t=1}^{\tau}\mathfrak R_t
        +\frac{3\cdot2^{13}\eta}\eps
        \sum_{t=1}^{\tau}\tr(H_t\Sigma_t).
        \label{eq:ctr-regret-energy}
    \end{align}
\end{lemma}

\begin{proof}
    Write
    \begin{align}
        U_t=\widetilde Y_t\Sigma_t^{-1/2}Z_t,
        \qquad
        b_t=\frac2\eps s_t.
    \end{align}
    The inequality $(u+v)^\top A(u+v)\leq2u^\top Au+2v^\top Av$ and
    $\Sigma_{t+1}\preceq3\Sigma_t$ give
    \begin{align}
        \normm{U_t+b_t}{\Sigma_{t+1}}^2
         & \leq2\normm{U_t}{\Sigma_{t+1}}^2
        +2\normm{b_t}{\Sigma_{t+1}}^2
        \\
         & \leq6\widetilde Y_t^2\norm{Z_t}^2
        +\frac{24}{\eps^2}s_t^\top\Sigma_ts_t.
    \end{align}
    Applying \eqref{eq:ctr-normal-width} and summing proves
    \eqref{eq:ctr-full-gradient-energy}.

    By \cref{lem:ctr-radial-coercivity} and $v(\mu_t)\leq\rho$,
    \begin{align}
        \sum_{t=1}^{\tau}v(\mu_t)^2
        \leq\rho\sum_{t=1}^{\tau}v(\mu_t)
        \leq\frac{\rho\eps}{2}
        \sum_{t=1}^{\tau}\mathfrak R_t.
    \end{align}
    Substituting this into \eqref{eq:ctr-full-gradient-energy} and multiplying
    by $\eta^2/2$ and $\eta/2$, respectively, proves
    \eqref{eq:ctr-focus-energy} and \eqref{eq:ctr-regret-energy}.
\end{proof}

\subsubsection{Centred concentration through the stopped interval}

\begin{proposition}[Centred concentration]
    \label{prop:ctr-concentration}
    Suppose the assumptions of \cref{lem:ctr-feedback} and
    \cref{lem:ctr-score-localisation} hold.  Suppose also that
    \begin{align}
        B_Yd^{3/2}\sqrt\ell
        \leq2^{-20}\sqrt n\,R_\rho,
        \label{eq:ctr-H-Freedman-scale}
    \end{align}
    and that $b\leq2^{-20}$ and $C_0\geq2^{12}$.  Then, with probability
    at least $1-\delta/4$,
    simultaneously,
    \begin{align}
        \left|
        \sum_{t=1}^{\tau}
        \left(
        \widehat q_t^{\mathrm{ctr}}(x_\star)-q_t(x_\star)
        \right)
        \right|
         & \leq
        \frac{2^{18}\sqrt b\,R_\rho\sqrt{n\ell}}\lambda,
        \label{eq:ctr-q-concentration}
        \\
        \sum_{t=1}^{\tau}
        \widetilde Y_t^2\norm{Z_t}^2
         & \leq2^{16}dn\ell R_\rho^2,
        \label{eq:ctr-random-gradient-energy}
    \end{align}
    and
    \begin{align}
         & -2^{22}R_\rho\ell^{3/2}\sqrt{dn}\,
        \overline\Sigma_{\tau+1}^{-1}
        \nonumber                             \\
         & \qquad\preceq
        \sum_{t=1}^{\tau}
        \left(\widehat H_t^{\mathrm{ctr}}-H_t\right)
        \nonumber                             \\
         & \qquad\preceq
        2^{22}R_\rho\ell^{3/2}\sqrt{dn}\,
        \overline\Sigma_{\tau+1}^{-1}.
        \label{eq:ctr-H-concentration}
    \end{align}
\end{proposition}

\begin{proof}
    We prove the three estimates separately, keeping all stopping indicators
    explicit.

    \proofstep{1}{The quadratic martingale}
    For $t\leq\tau$, put
    \begin{align}
        R_t=\normm{x_\star-\mu_t}{\Sigma_t^{-1}},
        \qquad
        h_t=
        \begin{cases}
            \Sigma_t^{-1/2}(x_\star-\mu_t)/R_t, & R_t>0, \\
            e_1,                                & R_t=0,
        \end{cases}
        \qquad
        G_t=\ip{h_t}{Z_t}.
    \end{align}
    Then $G_t$ is conditionally standard Gaussian and
    \begin{align}
        \widehat q_t^{\mathrm{ctr}}(x_\star)
        =\widetilde Y_tL_t+\kappa_t,
        \label{eq:ctr-q-score-representation}
    \end{align}
    where
    \begin{align}
        L_t=R_tG_t+\frac{\lambda R_t^2}{2}(G_t^2-1),
        \qquad
        \kappa_t=\frac2\eps\ip{s_t}{x_\star-\mu_t}.
    \end{align}
    The variable $\kappa_t$ is $\F_{t-1}$-measurable.  Hence
    \begin{align}
        \widehat q_t^{\mathrm{ctr}}(x_\star)-q_t(x_\star)
        =\widetilde Y_tL_t-\E_t[\widetilde Y_tL_t].
        \label{eq:ctr-q-martingale-cancellation}
    \end{align}
    The focus stop gives
    \begin{align}
        R_t\leq\sqrt{2B}
        =\frac{\sqrt{2b}}{\lambda\sqrt\ell}.
        \label{eq:ctr-R-focus}
    \end{align}
    Therefore
    \begin{align}
        L_t^2
         & \leq3R_t^2G_t^2
        +\frac{3\lambda^2R_t^4}{4}(G_t^4+1)
        \\
         & \leq
        \frac{12b}{\lambda^2}
        \left(
        \frac12+\frac{3}{2\ell^2}(G_t^2+1)^2
        \right).
        \label{eq:ctr-L-square}
    \end{align}
    Combining \eqref{eq:ctr-L-square} with
    \eqref{eq:ctr-score-localisation} gives
    \begin{align}
        \E_t[\widetilde Y_t^2L_t^2]
        \leq\frac{2^{20}b}{\lambda^2}
        \left(\E_t[\widetilde Y_t^2]+1\right).
        \label{eq:ctr-q-predictable-second}
    \end{align}

    Let
    \begin{align}
        \mathcal E_G
        =\left\{\max_{t\leq\tau}|G_t|\leq\sqrt\ell\right\}.
    \end{align}
    Since $G_t$ is conditionally standard Gaussian and
    $\1\{t\leq\tau\}$ is predictable,
    \begin{align}
        \Pp(\mathcal E_G^c)
        \leq2ne^{-\ell/2}
        \leq\frac\delta{64}.
        \label{eq:ctr-EG-probability}
    \end{align}
    On $\mathcal E_G$, \eqref{eq:ctr-R-focus} gives
    \begin{align}
        |L_t|
        \leq\frac{4\sqrt b}{\lambda}.
        \label{eq:ctr-L-maximum}
    \end{align}
    Let $\mathcal E_Y$ be the event in \cref{lem:ctr-feedback}.  Define
    \begin{align}
        Q_t=\1\{t\leq\tau\}\widetilde Y_tL_t,
        \qquad
        D_t=Q_t-\E_t[Q_t].
    \end{align}
    On $\mathcal E_G\cap\mathcal E_Y$,
    \begin{align}
        \sum_{t=1}^n\left(D_t^2+\E_t[D_t^2]\right)
         & \leq2\sum_{t=1}^n
        \left(Q_t^2+\E_t[Q_t^2]\right)
        \\
         & \leq\frac{2^{35}bnR_\rho^2}{\lambda^2},
        \label{eq:ctr-q-self-normaliser}
    \end{align}
    where the realised term uses
    \eqref{eq:ctr-L-maximum} and \eqref{eq:ctr-realised-Y2}, while the
    predictable term uses \eqref{eq:ctr-q-predictable-second} and
    \eqref{eq:ctr-predictable-Y2}.

    Apply the manuscript's self-normalised martingale inequality with
    \begin{align}
        z=\frac{2^{18}\sqrt b\,R_\rho\sqrt{n\ell}}\lambda,
        \qquad
        v=\frac{2^{35}bnR_\rho^2}{\lambda^2}.
    \end{align}
    Since $b\leq2^{-20}$,
    \begin{align}
        2\exp\left(-\frac{z^2}{2v}\right)
        \leq2e^{-\ell}
        \leq\frac\delta{64}.
    \end{align}
    Combining this with \eqref{eq:ctr-EG-probability} and the failure
    probability in \cref{lem:ctr-feedback} proves
    \eqref{eq:ctr-q-concentration} outside an event of probability at most
    $3\delta/32$.

    \proofstep{2}{The random part of the gradient energy}
    Let
    \begin{align}
        \mathcal E_Z
        =\left\{
        \max_{t\leq\tau}\norm{Z_t}^2
        \leq3(d+\ell)
        \right\}.
    \end{align}
    The chi-square upper-tail lemma and a predictable union bound give
    \begin{align}
        \Pp(\mathcal E_Z^c)\leq ne^{-\ell}\leq\frac\delta{64}.
        \label{eq:ctr-EZ-probability}
    \end{align}
    On $\mathcal E_Z\cap\mathcal E_Y$, using
    $d+\ell\leq2d\ell$ for $d,\ell\geq1$ gives
    \begin{align}
        \sum_{t=1}^{\tau}\widetilde Y_t^2\norm{Z_t}^2
         & \leq3(d+\ell)\sum_{t=1}^{\tau}\widetilde Y_t^2
        \\
         & \leq2^{16}dn\ell R_\rho^2.
    \end{align}
    This proves \eqref{eq:ctr-random-gradient-energy}.

    \proofstep{3}{The Hessian martingale}
    Let $\mathcal U$ be a $(2^{14}dn^2)^{-1}$-net of the Euclidean unit
    sphere.  The volumetric covering bound gives
    \begin{align}
        \log|\mathcal U|
        \leq d\log(1+2^{15}dn^2)
        \leq\frac{2d\ell}{C_0}.
        \label{eq:ctr-net-size}
    \end{align}
    Fix $u\in\mathcal U$ and define
    \begin{align}
        a_{t,u}=u^\top\Sigma_t^{-1}u,
        \qquad
        G_{t,u}
        =\ip{
            \frac{\Sigma_t^{-1/2}u}{\sqrt{a_{t,u}}}
        }{Z_t}.
        \label{eq:ctr-H-standardisation}
    \end{align}
    Conditionally, $G_{t,u}\sim\N(0,1)$, and
    \begin{align}
        Q_{t,u}
        =u^\top\widehat H_t^{\mathrm{ctr}}u
        =a_{t,u}\widetilde Y_t(G_{t,u}^2-1).
        \label{eq:ctr-H-direction}
    \end{align}

    Define the truncated score energy
    \begin{align}
        W_{t,u}
         & =\1\{t\leq\tau\}
        \widetilde Y_t^2(G_{t,u}^2-1)^2
        \1\{|\widetilde Y_t|\leq B_Y\}
        \1\{\norm{Z_t}^2\leq3(d+\ell)\}.
        \label{eq:ctr-W-definition}
    \end{align}
    By \cref{lem:ctr-score-localisation} and
    \eqref{eq:ctr-predictable-Y2},
    \begin{align}
        \sum_{t=1}^n\E_t[W_{t,u}]
         & \leq
        \sum_{t=1}^{\tau}
        \E_t\left[
            \widetilde Y_t^2(G_{t,u}^2-1)^2
            \right]
        \\
         & \leq2^{27}nR_\rho^2\ell^2.
        \label{eq:ctr-W-predictable}
    \end{align}
    Since $G_{t,u}^2\leq\norm{Z_t}^2$ on the truncation event,
    \begin{align}
        0\leq W_{t,u}
        \leq2^7B_Y^2d^2\ell^2
        \triangleq\Lambda.
        \label{eq:ctr-W-bound}
    \end{align}
    Put $R_{t,u}=W_{t,u}-\E_t[W_{t,u}]$.  Then
    \begin{align}
        \sum_{t=1}^n\E_t[R_{t,u}^2]
        \leq\Lambda\sum_{t=1}^n\E_t[W_{t,u}]
        \leq2^{34}B_Y^2d^2nR_\rho^2\ell^4.
        \label{eq:ctr-W-variance}
    \end{align}
    Apply Freedman's inequality with
    \begin{align}
        b=\Lambda,
        \qquad
        s=2^{17}B_Yd\sqrt n\,R_\rho\ell^2,
        \qquad
        z=\sqrt{2\log\frac{128|\mathcal U|}{\delta}}
        \leq\sqrt{\frac{d\ell}{8}}.
    \end{align}
    Indeed, \eqref{eq:ctr-net-size} and the definition of $\ell$ give
    \begin{align}
        \log\frac{128|\mathcal U|}{\delta}
        \leq\frac{3d\ell}{C_0}
        \leq\frac{d\ell}{16}.
    \end{align}
    Simultaneously for every
    $u\in\mathcal U$, outside an event of probability $\delta/64$,
    \begin{align}
        \sum_{t=1}^nR_{t,u}
         & \leq sz+\frac{bz^2}{3}
        \\
         & \leq
        2^{16}B_Yd^{3/2}\sqrt n\,R_\rho\ell^{5/2}
        +2^4B_Y^2d^3\ell^3
        \\
         & \leq2^{26}nR_\rho^2\ell^2.
        \label{eq:ctr-W-Freedman}
    \end{align}
    The last inequality follows from
    \eqref{eq:ctr-H-Freedman-scale}; squaring that condition also controls
    the second term.  Combining
    \eqref{eq:ctr-W-predictable} and \eqref{eq:ctr-W-Freedman} gives
    \begin{align}
        \sum_{t=1}^{\tau}
        \widetilde Y_t^2(G_{t,u}^2-1)^2
        \leq2^{28}nR_\rho^2\ell^2
        \label{eq:ctr-H-realised-directional-energy}
    \end{align}
    simultaneously for all $u\in\mathcal U$ on
    $\mathcal E_Y\cap\mathcal E_Z$.

    Let
    \begin{align}
        a_u^\star=\max_{t\leq\tau}a_{t,u}.
    \end{align}
    By \eqref{eq:ctr-H-direction},
    \eqref{eq:ctr-H-realised-directional-energy}, and
    \eqref{eq:ctr-W-predictable},
    \begin{align}
        \sum_{t=1}^{\tau}Q_{t,u}^2
         & \leq2^{28}(a_u^\star)^2nR_\rho^2\ell^2,
        \\
        \sum_{t=1}^{\tau}\E_t[Q_{t,u}^2]
         & \leq2^{27}(a_u^\star)^2nR_\rho^2\ell^2.
    \end{align}
    Hence, with
    \begin{align}
        D_{t,u}=\1\{t\leq\tau\}
        \left(Q_{t,u}-\E_t[Q_{t,u}]\right),
    \end{align}
    one has
    \begin{align}
        \sum_{t=1}^n
        \left(D_{t,u}^2+\E_t[D_{t,u}^2]\right)
        \leq2^{30}(a_u^\star)^2nR_\rho^2\ell^2.
        \label{eq:ctr-H-self-normaliser}
    \end{align}

    We next bound the range of $a_u^\star$.  The spectrum estimate in
    \cref{lem:ctr-trace} gives
    \begin{align}
        \norm{\overline\Sigma_{t+1}^{-1}}_{\op}^{1/2}
        \leq\sigma^{-1}+t\lambda\eta\Lip(e)
        \leq\sigma^{-1}+t,
        \label{eq:ctr-predictable-spectrum}
    \end{align}
    where the last inequality follows from the parameter conditions verified
    below.  Consequently,
    \begin{align}
        \norm{\overline\Sigma_{\tau+1}^{-1}}_{\op}
        \leq(\sqrt{2d}+n)^2
        \leq6dn^2.
        \label{eq:ctr-spectrum-upper}
    \end{align}
    Covariance comparison gives
    \begin{align}
        1\leq a_u^\star\leq9dn^2.
        \label{eq:ctr-a-range}
    \end{align}
    Let $J=\lceil\log_2(9dn^2)\rceil$.  For
    $j=0,\ldots,J$, put $b_j=2^j$.  On the event
    $b_{j-1}<a_u^\star\leq b_j$, the right-hand side of
    \eqref{eq:ctr-H-self-normaliser} is at most
    $2^{32}b_j^2nR_\rho^2\ell^2$.  The self-normalised martingale inequality
    therefore gives
    \begin{align}
         & \Pp\left(
        \left|\sum_{t=1}^nD_{t,u}\right|
        >2^{18}b_jR_\rho\ell
        \sqrt{n\log\frac{256|\mathcal U|(J+1)}\delta},
        \right.
        \nonumber             \\[-1mm]
         & \hspace{4cm}\left.
        \sum_{t=1}^n
        \left(D_{t,u}^2+\E_t[D_{t,u}^2]\right)
        \leq2^{32}b_j^2nR_\rho^2\ell^2
        \right)
        \leq\frac{\delta}{128|\mathcal U|(J+1)}.
    \end{align}
    Since $J+1\leq2^{4}dn^2$, \eqref{eq:ctr-net-size} and the definition
    of $\ell$ also give
    \begin{align}
        \log\frac{256|\mathcal U|(J+1)}\delta
        \leq\frac{6d\ell}{C_0}
        \leq\frac{d\ell}{16}.
    \end{align}
    A union bound over $u$ and $j$ therefore yields, outside an event of
    probability $\delta/128$,
    \begin{align}
        \left|
        u^\top\sum_{t=1}^{\tau}
        (\widehat H_t^{\mathrm{ctr}}-H_t)u
        \right|
        \leq2^{18}R_\rho\ell^{3/2}\sqrt{dn}
        \max_{t\leq\tau}u^\top\Sigma_t^{-1}u
        \label{eq:ctr-H-net}
    \end{align}
    for every $u\in\mathcal U$.

    Put
    \begin{align}
        M_\tau=\sum_{t=1}^{\tau}
        (\widehat H_t^{\mathrm{ctr}}-H_t),
        \qquad
        \overline A=\overline\Sigma_{\tau+1}^{-1}.
    \end{align}
    By covariance comparison and monotonicity of the predictable precision,
    \begin{align}
        \max_{t\leq\tau}\Sigma_t^{-1}
        \preceq\frac32\overline A.
    \end{align}
    Thus \eqref{eq:ctr-H-net} gives
    \begin{align}
        |u^\top M_\tau u|
        \leq3\cdot2^{17}R_\rho\ell^{3/2}\sqrt{dn}\,
        u^\top\overline A u.
        \label{eq:ctr-H-net-A}
    \end{align}
    Moreover, \eqref{eq:ctr-spectrum-upper} and
    $\overline A\succeq\sigma^{-2}I_d$ imply
    \begin{align}
        \frac{\lambda_{\max}(\overline A)}
        {\lambda_{\min}(\overline A)}
        \leq6dn^2.
        \label{eq:ctr-condition-number}
    \end{align}

    Fix a Euclidean unit vector $v$ and choose $u\in\mathcal U$ with
    $\norm{u-v}\leq(2^{14}dn^2)^{-1}$.  Symmetry gives
    \begin{align}
        |v^\top M_\tau v-u^\top M_\tau u|
        \leq2\norm{u-v}\norm{M_\tau}_{\op}
        \leq\frac{\norm{M_\tau}_{\op}}{2^{13}dn^2}.
        \label{eq:ctr-M-net-error}
    \end{align}
    Taking a supremum in the combination of
    \eqref{eq:ctr-H-net-A} and \eqref{eq:ctr-M-net-error} first yields
    \begin{align}
        \norm{M_\tau}_{\op}
        \leq2^{20}R_\rho\ell^{3/2}\sqrt{dn}\,
        \norm{\overline A}_{\op}.
        \label{eq:ctr-M-op-preliminary}
    \end{align}
    Also,
    \begin{align}
        |u^\top\overline A u-v^\top\overline A v|
         & \leq2\norm{u-v}\norm{\overline A}_{\op}
        \\
         & \leq\frac{12}{2^{14}}
        v^\top\overline A v,
        \label{eq:ctr-A-net-error}
    \end{align}
    where \eqref{eq:ctr-condition-number} is used in the last step.  Substituting
    \eqref{eq:ctr-M-op-preliminary} and \eqref{eq:ctr-A-net-error} into
    \eqref{eq:ctr-H-net-A}--\eqref{eq:ctr-M-net-error} gives
    \begin{align}
        |v^\top M_\tau v|
        \leq2^{22}R_\rho\ell^{3/2}\sqrt{dn}\,
        v^\top\overline A v.
    \end{align}
    Since this holds for every unit $v$, it is equivalent to
    \eqref{eq:ctr-H-concentration}.  The failure probabilities accumulated in
    the three steps are smaller than $\delta/4$.
\end{proof}

\subsubsection{Unconditional stability and regret}

The following deliberately non-optimised constants make every numerical
inequality in the proof explicit and mutually consistent:
\begin{align}
    C_0=2^{40},
    \qquad
    c_\lambda=2^{-80},
    \qquad
    c_\eta=2^{-80},
    \qquad
    b=2^{-20},
    \qquad
    C_\eps=2^{200}.
    \label{eq:ctr-concrete-constants}
\end{align}
Smaller $c_\lambda,c_\eta$ and larger $C_0,C_\eps$ also work.  No attempt is
made to optimise these universal constants.

\begin{theorem}[Tangent-centred unprojected regret]
    \label{thm:ctr-main}
    Let $\delta\in(0,1/2)$ and define
    \begin{align}
        \ell=C_0\log\frac{C_0nd}{\delta},
        \qquad
        \sigma=\frac1{M\sqrt{2d}},
        \qquad
        \lambda=\frac{c_\lambda}{d^{3/2}M\ell^2},
        \label{eq:ctr-main-parameters-a}
        \\
        \eta=\frac{c_\eta M}{\sqrt n\,\ell},
        \qquad
        \eps_0=\frac{C_\eps d^2M\ell^4}{\sqrt n},
        \label{eq:ctr-main-parameters-b}
    \end{align}
    with the constants in \eqref{eq:ctr-concrete-constants}.  If
    $\eps_0\leq1/2$, run \cref{alg:ctr-unprojected} with $\eps=\eps_0$.
    Otherwise play any fixed point in $\K$.  Then, with probability at least
    $1-\delta$,
    \begin{align}
        \operatorname{Reg}_n
        \leq2^{500}d^3M\ell^4\sqrt n.
        \label{eq:ctr-main-regret}
    \end{align}
    On the same event, the algorithm never fails and, on every round,
    \begin{align}
        v(\mu_t)
        \leq\rho
        =\frac{2\sqrt{2b}}{c_\lambda}d\ell^{3/2}.
        \label{eq:ctr-main-confinement}
    \end{align}
\end{theorem}

\begin{proof}
    Suppose first that $\eps_0>1/2$.  Then
    \begin{align}
        \sqrt n<2C_\eps d^2M\ell^4,
    \end{align}
    and therefore
    \begin{align}
        \operatorname{Reg}_n
        \leq n
        <2C_\eps d^2M\ell^4\sqrt n
        \leq2^{500}d^3M\ell^4\sqrt n.
    \end{align}
    Hence assume for the remainder that $\eps=\eps_0\leq1/2$.  This implies
    \begin{align}
        \sqrt n
        \geq2C_\eps d^2M\ell^4
        \geq2C_\eps d^{3/2}\ell^4,
        \label{eq:ctr-large-horizon-consequence}
    \end{align}
    because $M\geq d^{-1/2}$.  The constants in
    \eqref{eq:ctr-concrete-constants} also give $\rho\geq1$, and hence
    \begin{align}
        R_\rho=1+\rho\leq2\rho.
        \label{eq:ctr-R-rho-comparison}
    \end{align}

    \proofstep{1}{Verification of the stopped stochastic estimates}
    The predictable trace bound is independent of the sampled estimator.  Indeed,
    $\Lip(e)\leq8/\eps$, and direct substitution gives
    \begin{align}
        \lambda\eta\Lip(e)
         & \leq
        \frac{8c_\lambda c_\eta}
        {C_\eps d^{7/2}M\ell^7}
        \leq1.
        \label{eq:ctr-trace-small-step}
    \end{align}
    Since $\sigma\leq1$, the two hypotheses of
    \cref{lem:ctr-trace} follow from \eqref{eq:ctr-trace-small-step}.
    Therefore
    \begin{align}
        \eta\sum_{t=1}^{\tau_{\mathrm{cov}}}
        \tr(H_t\overline\Sigma_t)
        \leq\frac{4d\log(1+n)}\lambda.
        \label{eq:ctr-trace-budget}
    \end{align}
    In particular, by \eqref{eq:ctr-cov-comparison},
    \begin{align}
        \sum_{t=1}^{\tau}\tr(H_t\Sigma_t)
        \leq\frac{8d\log(1+n)}{\lambda\eta}.
        \label{eq:ctr-stopped-trace-budget}
    \end{align}

    We next verify the assumptions of \cref{lem:ctr-feedback}.  Since
    $R_\rho\geq\rho$ and
    \begin{align}
        \rho^2
        =\frac{8b}{c_\lambda^2}d^2\ell^3,
        \label{eq:ctr-rho-square}
    \end{align}
    the ratio between the left-hand side of
    \eqref{eq:ctr-feedback-param-a} and $nR_\rho^2$ is at most
    \begin{align}
        \frac{2^{12}d\ell^{3/2}}
        {\lambda\eta\eps\,n\rho^2}
         & \leq
        \frac{2^9c_\lambda}
        {b c_\eta C_\eps d^{3/2}M\ell^{5/2}}
        \leq2^{-100}.
        \label{eq:ctr-verify-feedback-a}
    \end{align}
    Also $e^{-8\ell}/\eps^2\leq R_\rho^2$, because
    $\eps^{-2}\leq n$ and $ne^{-8\ell}\leq1$.  Furthermore,
    \begin{align}
        \frac{B_Y\sqrt\ell}{\sqrt nR_\rho}
         & =\frac{2^8\ell^{3/2}}{\sqrt n}
        \left(1+\frac1{\eps R_\rho}\right)
        \\
         & \leq
        \frac{2^7}{C_\eps d^{3/2}\ell^{5/2}}
        +\frac{2^8\ell^{3/2}}
        {C_\eps d^2M\ell^4R_\rho}
        \leq2^{-100}.
        \label{eq:ctr-verify-feedback-c}
    \end{align}
    Thus \eqref{eq:ctr-feedback-param-a}--\eqref{eq:ctr-feedback-param-d}
    hold.  Similarly,
    \begin{align}
        2^{20}\eps^{-2}R_\rho^2
         & \leq
        \frac{2^{25}bn}
        {c_\lambda^2C_\eps^2d^2M^2\ell^5}
        \leq2^{25}n
        \leq e^{4\ell},
        \label{eq:ctr-verify-score-tail}
    \end{align}
    proving \eqref{eq:ctr-score-tail-parameter}.  Finally,
    \begin{align}
        \frac{B_Yd^{3/2}\sqrt\ell}{\sqrt nR_\rho}
         & =\frac{2^8d^{3/2}\ell^{3/2}}{\sqrt n}
        \left(1+\frac1{\eps R_\rho}\right)
        \leq2^{-100},
        \label{eq:ctr-verify-H-Freedman}
    \end{align}
    by \eqref{eq:ctr-large-horizon-consequence}.  Hence
    \eqref{eq:ctr-H-Freedman-scale} holds.  The net logarithm condition follows
    from $C_0=2^{40}$.  All hypotheses of
    \cref{prop:ctr-concentration} are therefore satisfied without assuming a
    feedback good event.

    Work on the event in \cref{prop:ctr-concentration}.  It has probability at
    least $1-\delta/4$.

    \proofstep{2}{Covariance stability}
    From \eqref{eq:ctr-H-concentration},
    \begin{align}
        2^{22}\lambda\eta R_\rho\ell^{3/2}\sqrt{dn}
         & \leq
        2^{23}\lambda\eta\rho\ell^{3/2}\sqrt{dn}
        \\
         & =2^{25}\eta\sigma\sqrt b\,\ell\sqrt{dn}
        =\frac{2^{25}}{\sqrt2}c_\eta\sqrt b
        \leq\frac14.
        \label{eq:ctr-H-small}
    \end{align}
    Multiplying \eqref{eq:ctr-H-concentration} by $\lambda\eta$ therefore
    gives
    \begin{align}
        -\frac14\overline\Sigma_{\tau+1}^{-1}
        \preceq
        \lambda\eta\sum_{t=1}^{\tau}
        (\widehat H_t^{\mathrm{ctr}}-H_t)
        \preceq
        \frac14\overline\Sigma_{\tau+1}^{-1}.
        \label{eq:ctr-H-small-matrix}
    \end{align}
    All earlier covariance updates were performed, so the proposed precision
    at round $\tau$ is
    \begin{align}
        \Sigma_\tau^{-1}+\lambda\eta\widehat H_\tau^{\mathrm{ctr}}
         & =\Sigma_1^{-1}
        +\lambda\eta\sum_{t=1}^{\tau}
        \widehat H_t^{\mathrm{ctr}}
        \\
         & =\overline\Sigma_{\tau+1}^{-1}
        +\lambda\eta\sum_{t=1}^{\tau}
        (\widehat H_t^{\mathrm{ctr}}-H_t).
    \end{align}
    Hence
    \begin{align}
        \frac34\overline\Sigma_{\tau+1}^{-1}
        \preceq
        \Sigma_\tau^{-1}+\lambda\eta\widehat H_\tau^{\mathrm{ctr}}
        \preceq
        \frac54\overline\Sigma_{\tau+1}^{-1}.
        \label{eq:ctr-strict-covariance}
    \end{align}
    The update is therefore performed at round $\tau$, and covariance
    comparison cannot fail there.  Moreover, for every $t\leq\tau$,
    \begin{align}
        \Sigma_{t+1}^{-1}
        \succeq\frac12\overline\Sigma_{t+1}^{-1}
        \succeq\frac12\overline\Sigma_t^{-1}
        \succeq\frac13\Sigma_t^{-1},
    \end{align}
    so
    \begin{align}
        \Sigma_{t+1}\preceq3\Sigma_t.
        \label{eq:ctr-next-covariance}
    \end{align}

    \proofstep{3}{Focus stability and inward drift}
    Summing \eqref{eq:ctr-unprojected-ONS} with $x=x_\star$ gives
    \begin{align}
        \frac12
        \normm{\mu_{\tau+1}-x_\star}{\Sigma_{\tau+1}^{-1}}^2
         & =\frac12
        \normm{\mu_1-x_\star}{\Sigma_1^{-1}}^2
        +\eta\sum_{t=1}^{\tau}q_t(x_\star)
        \nonumber   \\
         & \quad
        +\eta\sum_{t=1}^{\tau}
        \left(
        \widehat q_t^{\mathrm{ctr}}(x_\star)-q_t(x_\star)
        \right)
        +\frac{\eta^2}{2}
        \sum_{t=1}^{\tau}
        \normm{\widehat g_t^{\mathrm{ctr}}}{\Sigma_{t+1}}^2.
        \label{eq:ctr-focus-decomposition}
    \end{align}

    Regret switching and the bracket proposition give
    \begin{align}
        q_t(x_\star)
        \leq
        -\mathfrak R_t
        +\frac{17}{16}\tr(H_t\Sigma_t)
        +\operatorname{tail}_t,
        \label{eq:ctr-q-inward}
    \end{align}
    where
    \begin{align}
        \operatorname{tail}_t
        =\frac{\lambda\Lip(e)e^{-\ell}}{\sqrt\pi}
        \norm{\mu_t-x_\star}
        \normm{\mu_t-x_\star}{\Sigma_t^{-1}}.
    \end{align}
    The bracket scale conditions hold because
    \begin{align}
        \lambda(d+2\ell)
        \leq\frac{4c_\lambda}{C_0}\leq\frac1{64},
        \qquad
        \lambda\sqrt{2B}\sqrt{2\ell}
        =2\sqrt b\leq\frac1{64}.
    \end{align}
    The covariance lower bound and focus give
    \begin{align}
        \norm{\mu_t-x_\star}\leq2\sigma\sqrt B,
        \qquad
        \normm{\mu_t-x_\star}{\Sigma_t^{-1}}\leq\sqrt{2B}.
    \end{align}
    Thus
    \begin{align}
        \sum_{t=1}^{\tau}\operatorname{tail}_t
        \leq
        \frac{32bn\sigma e^{-\ell}}{\lambda\eps\ell}
        \leq1.
        \label{eq:ctr-bracket-tail-sum}
    \end{align}
    The last inequality follows by substituting the parameters and using
    $n^{3/2}e^{-\ell}\leq1$.

    By \eqref{eq:ctr-next-covariance},
    \cref{lem:ctr-gradient-energy} applies.  Furthermore,
    \begin{align}
        \frac{24\eta\rho}{\eps}
        =\frac{48\sqrt{2b}\,c_\eta}
        {c_\lambda C_\eps d\ell^{7/2}}
        \leq\frac12.
        \label{eq:ctr-regret-absorption-condition}
    \end{align}
    Combining \eqref{eq:ctr-q-inward} with
    \eqref{eq:ctr-focus-energy} therefore absorbs the predictable tangent
    energy into the negative inward-drift term $-\eta\sum_t\mathfrak R_t$.
    Dropping the remaining nonpositive term and using
    \eqref{eq:ctr-q-concentration},
    \eqref{eq:ctr-random-gradient-energy}, and
    \eqref{eq:ctr-stopped-trace-budget} yields
    \begin{align}
        \frac12
        \normm{\mu_{\tau+1}-x_\star}{\Sigma_{\tau+1}^{-1}}^2
         & \leq4d^3M^2
        +\frac{2^{18}\sqrt b\,\eta R_\rho\sqrt{n\ell}}\lambda
        +3\cdot2^{16}\eta^2dn\ell R_\rho^2
        \nonumber      \\
         & \quad
        +\frac{17d\log(1+n)}{2\lambda}
        +\frac{3\cdot2^{16}\eta d\log(1+n)}{\lambda\eps}
        +\eta.
        \label{eq:ctr-focus-budget}
    \end{align}
    Each term is at most $B/8$ under
    \eqref{eq:ctr-concrete-constants}.  Explicitly, after division by
    $B=b/(\lambda^2\ell)$, the first four non-negligible ratios are bounded by
    \begin{align}
        \frac{4c_\lambda^2}{b\ell^3},
        \qquad
        \frac{2^{22}c_\eta}{\sqrt d},
        \qquad
        \frac{3\cdot2^{22}c_\eta^2}{\ell},
        \qquad
        \frac{17c_\lambda}{2b\sqrt d\,M},
    \end{align}
    respectively; each is smaller than $2^{-20}$.  The last two terms in
    \eqref{eq:ctr-focus-budget} are smaller still because they contain
    $C_\eps^{-1}$ or $B^{-1}$.  Therefore
    \begin{align}
        \frac12
        \normm{\mu_{\tau+1}-x_\star}{\Sigma_{\tau+1}^{-1}}^2
        \leq B.
        \label{eq:ctr-focus-strict}
    \end{align}
    Focus cannot fail at round $\tau$.

    If $\tau<n$, \eqref{eq:ctr-strict-covariance} and
    \eqref{eq:ctr-focus-strict} imply
    \begin{align}
        \tau<\tau_{\mathrm{cov}}\wedge\tau_{\mathrm{focus}}=\tau,
    \end{align}
    a contradiction.  Hence
    \begin{align}
        \tau=n.
        \label{eq:ctr-global-stability}
    \end{align}
    In particular, the algorithm never fails, and
    \eqref{eq:ctr-main-confinement} follows from
    \cref{lem:ctr-coarse-confinement}.

    \proofstep{4}{Extension regret}
    Continue on the event used above.  Regret switching and the bracket bound
    give
    \begin{align}
        \sum_{t=1}^n\mathfrak R_t
        \leq
        -\sum_{t=1}^nq_t(x_\star)
        +\frac{17}{16}\sum_{t=1}^n\tr(H_t\Sigma_t)+1.
        \label{eq:ctr-extension-switch}
    \end{align}
    By \eqref{eq:ctr-unprojected-ONS}, dropping the nonnegative final
    potential gives
    \begin{align}
        -\sum_{t=1}^n\widehat q_t^{\mathrm{ctr}}(x_\star)
        \leq
        \frac{4d^3M^2}{\eta}
        +\frac\eta2\sum_{t=1}^n
        \normm{\widehat g_t^{\mathrm{ctr}}}{\Sigma_{t+1}}^2.
        \label{eq:ctr-ONS-regret}
    \end{align}
    Substitute
    $-q_t=-\widehat q_t^{\mathrm{ctr}}
        +(\widehat q_t^{\mathrm{ctr}}-q_t)$ into
    \eqref{eq:ctr-extension-switch}, and apply
    \eqref{eq:ctr-regret-energy}.  The absorption condition
    \eqref{eq:ctr-regret-absorption-condition} moves the resulting multiple of
    $\sum_t\mathfrak R_t$ to the left.  Using
    \eqref{eq:ctr-q-concentration},
    \eqref{eq:ctr-random-gradient-energy}, and
    \eqref{eq:ctr-stopped-trace-budget} gives
    \begin{align}
        \frac12\sum_{t=1}^n\mathfrak R_t
         & \leq
        \frac{4d^3M^2}{\eta}
        +3\cdot2^{16}\eta dn\ell R_\rho^2
        +\frac{2^{18}\sqrt b\,R_\rho\sqrt{n\ell}}\lambda
        \nonumber \\
         & \quad
        +\frac{17d\log(1+n)}{2\lambda\eta}
        +\frac{3\cdot2^{16}d\log(1+n)}{\lambda\eps}
        +1.
        \label{eq:ctr-extension-budget}
    \end{align}
    Substituting \eqref{eq:ctr-main-parameters-a}--
    \eqref{eq:ctr-main-parameters-b}, using
    $R_\rho\leq2\rho$, \eqref{eq:ctr-rho},
    $M\sqrt d\geq1$, and $\log(1+n)\leq\ell$, yields
    \begin{align}
        \sum_{t=1}^n\mathfrak R_t
        \leq2^{400}d^3M\ell^4\sqrt n.
        \label{eq:ctr-extension-regret}
    \end{align}
    For example, the five nonconstant terms in
    \eqref{eq:ctr-extension-budget} are respectively of orders
    \begin{align}
        d^3M\ell\sqrt n,
        \quad
        d^3M\ell^3\sqrt n,
        \quad
        d^{5/2}M\ell^4\sqrt n,
        \quad
        d^{5/2}\ell^4\sqrt n,
        \quad
        \sqrt d\,\ell^{-1}\sqrt n,
    \end{align}
    each bounded by the right-hand side of
    \eqref{eq:ctr-extension-regret}.

    \proofstep{5}{Realised regret}
    Azuma's inequality gives, outside an event of probability $\delta/4$,
    \begin{align}
        \sum_{t=1}^n
        \left(f(A_t)-\E_t[f(A_t)]\right)
        \leq\sqrt{2n\ell}.
        \label{eq:ctr-action-martingale}
    \end{align}
    As in the manuscript,
    \begin{align}
        f(x_\star)-\inf_{x\in\K}f(x)\leq\eps,
        \qquad
        f(A_t)\leq e(X_t).
    \end{align}
    Combining these inequalities with
    \eqref{eq:ctr-extension-regret} and
    $n\eps=C_\eps d^2M\ell^4\sqrt n$ gives
    \begin{align}
        \operatorname{Reg}_n
         & \leq
        \sqrt{2n\ell}
        +2^{400}d^3M\ell^4\sqrt n
        +C_\eps d^2M\ell^4\sqrt n
        \\
         & \leq2^{500}d^3M\ell^4\sqrt n.
    \end{align}
    The concentration event and action-martingale event fail with total
    probability at most $\delta/2<\delta$, completing the proof.
\end{proof}

\subsubsection{Why the direct small-shell distance argument does not close}

The preceding theorem obtains the deterministic shell
$v(\mu_t)\leq\rho=O(d\ell^{3/2})$ from the focus potential.  The following
identity is the natural starting point for trying to replace this shell by
$v(\mu_t)=O(\sqrt\ell)$ while retaining the larger step size used in the
projected theorem.

\begin{lemma}[Exact stopped distance-to-$\K_\eps$ recursion]
    \label{lem:ctr-distance-recursion}
    On an active round, let
    \begin{align}
        P_t
        =\argmin_{x\in\K_\eps}
        \normm{\mu_t-x}{\Sigma_t^{-1}}^2,
        \qquad
        \mathsf D_t
        =\frac12\normm{\mu_t-P_t}{\Sigma_t^{-1}}^2.
    \end{align}
    The minimizer is unique and $\F_{t-1}$-measurable.  Put
    \begin{align}
        R_t=\sqrt{2\mathsf D_t},
        \qquad
        h_t=
        \begin{cases}
            \Sigma_t^{-1/2}(\mu_t-P_t)/R_t, & R_t>0, \\
            e_1,                            & R_t=0,
        \end{cases}
        \qquad
        G_t=\ip{h_t}{Z_t}.
    \end{align}
    Then
    \begin{align}
        \mathsf D_{t+1}-\mathsf D_t
         & \leq
        -\frac{2\eta}{\eps}v(\mu_t)
        +\eta\widetilde Y_t
        \left(
        -R_tG_t
        +\frac{\lambda R_t^2}{2}(G_t^2-1)
        \right)
        \nonumber \\
         & \qquad
        +\frac{\eta^2}{2}
        \normm{\widehat g_t^{\mathrm{ctr}}}{\Sigma_{t+1}}^2.
        \label{eq:ctr-distance-recursion}
    \end{align}
\end{lemma}

\begin{proof}
    Strict convexity of the quadratic objective gives uniqueness of $P_t$.
    To prove measurability, let
    \begin{align}
        F(\mu,A,x)=(x-\mu)^\top A(x-\mu),
        \qquad A\succ0.
    \end{align}
    The function $F$ is continuous, $\K_\eps$ is compact, and the minimizer is
    unique.  If $(\mu_k,A_k)\to(\mu,A)$ and $P(\mu_k,A_k)$ has a convergent
    subsequence with limit $p$, continuity and optimality imply that $p$ is
    the unique minimizer for $(\mu,A)$.  Hence the entire sequence converges
    to $P(\mu,A)$, so the projection map is continuous and therefore Borel
    measurable.

    By the definition of $\mathsf D_{t+1}$,
    \begin{align}
        \mathsf D_{t+1}
        \leq\frac12
        \normm{\mu_{t+1}-P_t}{\Sigma_{t+1}^{-1}}^2.
    \end{align}
    Apply \eqref{eq:ctr-unprojected-ONS} with $x=P_t$.  The random score part
    satisfies
    \begin{align}
         & \ip{
            \widetilde Y_t\Sigma_t^{-1/2}Z_t
        }{P_t-\mu_t}
        +\frac\lambda2
        (P_t-\mu_t)^\top\widehat H_t^{\mathrm{ctr}}
        (P_t-\mu_t)
        \\
         & \qquad=
        \widetilde Y_t
        \left(
        -R_tG_t
        +\frac{\lambda R_t^2}{2}(G_t^2-1)
        \right).
    \end{align}
    Finally, \eqref{eq:ctr-inward-support} gives
    \begin{align}
        \frac2\eps\ip{s_t}{P_t-\mu_t}
        \leq-\frac2\eps v(\mu_t).
    \end{align}
    Substitution proves \eqref{eq:ctr-distance-recursion}.
\end{proof}

The threshold implication itself is valid:
\begin{align}
    v(\mu_t)>r
    \quad\Longrightarrow\quad
    \mathsf D_t>\frac{r^2}{16\sigma^2},
    \label{eq:ctr-valid-distance-threshold}
\end{align}
by the proof of \cref{lem:ctr-coarse-confinement}.  What fails is the
increment inequality needed to telescope
\eqref{eq:ctr-distance-recursion} over an excursion.  If
\begin{align}
    \nu_r
    =\max\{t<\tau_r:v(\mu_t)\leq r/2\},
    \qquad
    \tau_r=\min\{t:v(\mu_t)>r\},
\end{align}
then the direct excursion proof would require a deterministic lower bound such
as
\begin{align}
    \mathsf D_{\tau_r}-\mathsf D_{\nu_r}
    \geq c\frac{r^2}{\sigma^2}
    \label{eq:ctr-false-excursion-increment}
\end{align}
for a universal $c>0$.  This is false, even with a fixed covariance.

\begin{proposition}[Failure of the distance-excursion increment]
    \label{prop:ctr-distance-excursion-failure}
    For every $r>0$ and every $c>0$, there are a convex body, a fixed
    positive-definite covariance, and points $x,y$ such that
    \begin{align}
        v(x)\leq\frac r2,
        \qquad
        v(y)>r,
        \qquad
        \frac12\operatorname{dist}_{\Sigma^{-1}}(y,\K_\eps)^2
        -\frac12\operatorname{dist}_{\Sigma^{-1}}(x,\K_\eps)^2
        <-c\frac{r^2}{\norm{\Sigma}_{\op}}.
        \label{eq:ctr-distance-excursion-counterexample}
    \end{align}
    Moreover, Euclidean squared distance cannot replace the Mahalanobis
    potential using only the deterministic correction $-(2\eta/\eps)\Sigma s$
    for an arbitrary selector $s\in\partial v(\mu)$.
\end{proposition}

\begin{proof}
    Take
    \begin{align}
        \K_\eps=[-1,1]^2,
        \qquad
        v(z)=\left(\max\{|z_1|,|z_2|\}-1\right)_+.
    \end{align}
    Put
    \begin{align}
        x=(1+r/2,0),
        \qquad
        y=(0,1+2r),
        \qquad
        \Sigma^{-1}=
        \begin{pmatrix}
            \kappa & 0 \\
            0      & 1
        \end{pmatrix}.
    \end{align}
    Then $v(x)=r/2$ and $v(y)=2r>r$, while
    \begin{align}
        \frac12\operatorname{dist}_{\Sigma^{-1}}(x,\K_\eps)^2
        =\frac{\kappa r^2}{8},
        \qquad
        \frac12\operatorname{dist}_{\Sigma^{-1}}(y,\K_\eps)^2
        =2r^2.
    \end{align}
    Since $\norm{\Sigma}_{\op}=1$, choosing
    $\kappa>16+8c$ proves
    \eqref{eq:ctr-distance-excursion-counterexample}.  Thus a larger raw
    violation can be reached while the stopped Mahalanobis distance decreases
    by an arbitrarily large amount.

    For the Euclidean claim, take
    \begin{align}
        \mu=(1+r,1+r),
        \qquad
        s=e_1\in\partial v(\mu),
        \qquad
        \Sigma=
        \begin{pmatrix}
            1  & -2 \\
            -2 & 5
        \end{pmatrix}.
    \end{align}
    The matrix is positive definite because its leading principal minors are
    $1$ and $1$.  The Euclidean projection is $P=(1,1)$, and for the pure
    tangent correction $\Delta=-a\Sigma s=a(-1,2)$ with $a>0$,
    \begin{align}
        \ip{\mu-P}{\Delta}=ar>0.
    \end{align}
    Hence
    \begin{align}
        \frac12\operatorname{dist}(\mu+\Delta,\K_\eps)^2
        -\frac12\operatorname{dist}(\mu,\K_\eps)^2
        =ar+O(a^2)>0
    \end{align}
    for all sufficiently small $a$.  This disproves a selector-uniform
    Euclidean inward-drift inequality based only on the deterministic tangent
    term.  It does not assert that the full centred estimator is unstable; the
    random tangent remainder has a compensating conditional mean.
\end{proof}

\paragraph{Exact proof status for the stronger parameter choice.}
The centred estimator removes the affine gauge component from the random
Hermite scores, and every estimator, feedback, covariance, focus, and regret
calculation above is complete.  With the smaller step size in
\eqref{eq:ctr-main-parameters-b}, the stopped focus potential itself supplies
coarse confinement and closes the bootstrap.  To retain the projected
manuscript's larger choice $\eta\asymp dM/\sqrt n$, the covariance
concentration bound would require a shell $v(\mu_t)=O(\sqrt\ell)$ rather than
\eqref{eq:ctr-main-confinement}.  The direct attempt to prove that shell from
\eqref{eq:ctr-distance-recursion} first needs an excursion comparison of the
form \eqref{eq:ctr-false-excursion-increment};
\cref{prop:ctr-distance-excursion-failure} shows that this particular
inequality is false.  A same-rate theorem would therefore require a different
multi-face confinement potential or a more restrictive selector rule.  No
such extra assumption is used in \cref{thm:ctr-main}.
